\Needspace{10\baselineskip}
\pdfbookmark[0]{Anzeigen. Opérations géodésiques et astronomiques. 1830.}{anzeigen-operations-1830}
\begin{center}
\rule{.85\textwidth}{.3pt}\par
Göttingische gelehrte Anzeigen.\quad 27. Februar 1830.\par
\rule{.85\textwidth}{.3pt}
\end{center}
\hspace*{\parindent}\textit{Opérations géodésiques et astronomiques pour la mesure d’un arc du parallèle moyen, exécutées en Piémont et en Savoie par une commission composée d’officiers de l’état major général et d’astronomes Piémontais et Autrichiens en} 1821, 1822, 1823. \textit{Milan, de l’imprimerie impériale et royale: Tome premier} 1825. 238 Seiten. \textit{Tome second} 1827. 412 Seiten im Quartformat. Dazu ein Heft mit Figuren, Karten und sechs Rundansichten.

Die Idee zu der großen Längengradmessung, zu deren Hauptbestandteilen die in diesem Werk veröffentlichten Arbeiten gehören, ging ursprünglich von \textsc{Laplace} aus. Seit 1802 hatten französische Ingenieure in Oberitalien ausgedehnte Dreiecksmessungen ausgeführt, zunächst für topografisch-militärische
